\begin{proof}
(i) Table \ref{tab:c11-defectcounts} and the evaluation of the leg set: a single output $Y$ gives $\operatorname{tr}_{\mathfrak g}\mathrm{ad}_Y=0$ on the colour singlet; $s_0$ has the legs $A|_{W_{\mathbb R}^\perp}$, $\lambda^*$, and its bracket with a letter needs the leg $A^*|_{W_{\mathbb R}^\perp}$. Theorem \ref{thm:c11-oddsimplex} (ii) gives the position part. (ii) The leg set $A|_V$ has a single output; the empty leg set is homogeneous in $\varepsilon$ of degree $-\tfrac12$ (Lemma \ref{lem:c10-self}). Theorem \ref{thm:c11-oddsimplex} (iii) gives the two-body parts; by Corollary \ref{cor:c7-K} the term in the second derivatives of $v$ and the term in the antisymmetric part of $dv$ vanish on a scalar bilinear. The one-body part: $-\{-\tfrac{22}3\,\text{units}\int(b,x)(A,A^*),O\}=\tfrac{44}3(b,x)\,O$ units, the rotation and spin terms of the solution acting on $A|_V$ and $\lambda$ alone, and $-\theta^K_\perp$ units times $\{S^{(1)}(K_b),O\}$. With $v(x)^\mu\partial_\mu O$ and $\{S^{(1)}(K_b),O\}=v^\mu\partial_\mu O+\tfrac24(\operatorname{div}v)O$ (Definition \ref{def:c7-vdot}), $-\tfrac13v^\mu\partial_\mu O-\tfrac{40}3(b,x)O+\tfrac{44}3(b,x)O=-\tfrac13\{S^{(1)}(K_b),O\}$ on the traceless part, and on the trace $-(b,x)\Upsilon^0_\Delta$ adds $(\tfrac{104}3+\tfrac{40}3)(b,x)O=48(b,x)O$, the one-body part being the same. Proposition \ref{prop:c8-primary}: $K_1(O)=-2(b,x)\kappa\,O$ for $D_1(O)=2\kappa\,O$.
\end{proof}
